\clearpage
\chapter*{강연 195의 정오표}
\addcontentsline{toc}{chapter}{강연 195의 정오표}
\setcounter{footnote}{0}
\src{300}

\noindent\textbf{4. TDTE II: 형식적 모듈라이 이론의 존재 정리.}

\begin{center}
[부르바키 세미나, 제12권, 1959/60, 제195호, 22쪽.]
\end{center}

\noindent\textbf{8쪽.} --- 명제 5.1에 쓰인 공식은 $\Lambda$가 체인
경우에만 옳다. 일반적인 경우에는
$\mathfrak m_\xi/\mathfrak m_\xi^2$을 이 공간을
$\mathfrak n_\xi/\mathfrak n_\xi^2$의 상으로 나눈 몫으로 바꾸어야 한다.
여기서 $\mathfrak n$은 $\Lambda$의 극대 아이디얼이다. 또한 $O$가
$\Lambda$ 위에서 매끄럽다는 것에 관해 주어진 정의는 잉여체 확대 $k'/k$가
분리 가능할 때에만 옳다. 일반적인 경우에는 SGA III, 1.1을 보라.

\noindent\textbf{14쪽, 주.} --- 여기서 제기한 문제들은 사영의 경우
``힐베르트 스킴''에 의해 완전히 해결된다(cf. TDTE IV). 그러나 Nagata와
Hironaka의 예들은 사영성 가정을 두지 않으면 고려한 함자들이 반드시 표현
가능한 것은 아님을 보여 준다. 이는 차원 3인 비특이 완비 다양체의 0차원
부분다양체들을 분류하는 데에만 국한하더라도 마찬가지이다. 이 사실은 그러한
다양체의 2차 대칭곱이 존재하지 않을 수 있다는 점과 관련된다.

\src{301}
\noindent\textbf{15쪽, 제3항.} --- 더 자세한 연구는 TDTE V를 보라.

\noindent\textbf{16쪽, 명제 3.1.} --- ``$f$가 고유이면'' 대신
``$f$가 고유이고 분리 가능하면''으로 읽어야 한다.

\noindent\textbf{18쪽과 19쪽, 주, 제2항.} --- 나는 최근 체 위의 아벨
다양체에 대한 형식적 모듈라이 스킴이 실제로 비트 환 위에서 매끄럽다는 것을
증명했다. 다시 말해, 어떤 다른 환의 몫인 아르틴 국소환 위의 모든 아벨
스킴은 그 다른 환 위의 아벨 스킴을 환원하여 얻어진다. 증명은 강연 182,
12쪽에서 도입한 올림 장애류의 함자적 성질만을 사용한다. 또한 종수 $g$인
곡선들의 모듈라이 스킴과 편극된 아벨 스킴들의 모듈라이 스킴을 Mumford가
구성했음을 상기하자(cf. SHMT).

\noindent\textbf{19쪽, 제5문단.} --- $U$를 정의하는 데 쓰는 단면은
$\mathfrak X$ 위에서뿐 아니라 모든 $X_n$ 위에서도 영인자가 아니라고
가정해야 한다.

\noindent\textbf{20쪽, 12--13행.} ---
\[
  G_x=G\bigl(\operatorname{Spec}(\widehat{\mathcal O}_{\mathfrak X,x})\bigr),
\]
대신
\[
  G_x=G\bigl(\operatorname{Spec}(\mathcal O_{\mathfrak X,x})\bigr)
  \subset
  G\bigl(\operatorname{Spec}(\widehat{\mathcal O}_{\mathfrak X,x})\bigr),
\]
으로 읽고, ``$\mathcal O'_{0x}$'' 대신
``$\widehat{\mathcal O}'_{0x}$''으로 읽어야 한다.
